\section*{Hoofdstuk \ref{ch:g11:titration} --- Titratie}

\begin{solution}{exo:g11:titration:1}
De \omterm{def:g11:titration:titration}{titrant} is de oplossing met de nauwkeurig bekende concentratie, in de
buret; de \omterm{def:g11:titration:titration}{getitreerde oplossing} bevat het deeltje dat gemeten moet
worden, in de erlenmeyer. Bij het \omterm{def:g11:titration:equivalence}{equivalentiepunt} zijn de toegevoegde
\omterm{def:g11:titration:titration}{titrant} en het getitreerde deeltje in de verhouding van de
\omterm{def:g8:balanced-equations:equation}{reactievergelijking} aanwezig: beide zijn verbruikt.
\end{solution}

\begin{solution}{exo:g11:titration:2}
$n = 0.0200 \times 0.0125 = \qty{2.50e-4}{mol}$.
\end{solution}

\begin{solution}{exo:g11:titration:3}
$V_E = \qty{14.3}{mL}$; de lijn van ijzer(II) begint bij
\qty{14.3e-4}{mol}, dat is \qty{1.43e-3}{mol}.
\end{solution}

\begin{solution}{exo:g11:titration:4}
De \omterm{def:g11:titration:titration}{titrant} is zelf gekleurd (paars) en zijn product bijna kleurloos: de
kleur verschijnt in de erlenmeyer zodra het permanganaat in overmaat is.
\end{solution}

\begin{solution}{exo:g11:titration:5}
5: volgens de vergelijking worden er 5 ijzer(II)\omterm{def:g9:ions:ion}{ionen} verbruikt per
permanganaation, en bij het \omterm{def:g11:titration:equivalence}{equivalentiepunt} zijn de \omterm{def:g7:chemical-reactions:reactant}{beginstoffen} in
precies die verhouding bij elkaar gevoegd.
\end{solution}

\begin{solution}{exo:g11:titration:6}
$n(\ce{MnO4^-}) = 0.0200 \times 0.00840 = \qty{1.68e-4}{mol}$;
$n(\ce{Fe^{2+}}) = 5 \times \num{1.68e-4} = \qty{8.40e-4}{mol}$;
$c = \num{8.40e-4} / 0.0100 = \qty{0.0840}{mol/L}$;
$C_m = 0.0840 \times 55.8 = \qty{4.69}{g/L}$.
\end{solution}

\begin{solution}{exo:g11:titration:7}
Verhouding 1:1: $n = 0.0100 \times 0.0114 = \qty{1.14e-4}{mol}$ in
\qty{10.0}{mL}, dus \qty{1.14e-3}{mol} in de \qty{100.0}{mL}, dat is
$\num{1.14e-3} \times 176.0 = \qty{0.201}{g}$: ongeveer \qty{200}{mg}.
\end{solution}

\begin{solution}{exo:g11:titration:8}
Een druppel zou niet meteen reageren: de kleur van de \omterm{def:g11:titration:titration}{titrant} zou vóór
het \omterm{def:g11:titration:equivalence}{equivalentiepunt} blijven hangen, en je zou het einde te vroeg zien
(of je zou na elke druppel minutenlang moeten wachten).
\end{solution}

\begin{solution}{exo:g11:titration:9}
Te groot: ze is voorbij het \omterm{def:g11:titration:equivalence}{equivalentiepunt} gegaan en heeft
permanganaat in overmaat toegevoegd. Vlak bij het \omterm{def:g11:titration:equivalence}{equivalentiepunt} had ze
de \omterm{def:g11:titration:titration}{titrant} druppel voor druppel moeten toevoegen, onder zwenken, en
moeten stoppen bij de eerste blijvende zwakke roze kleur.
\end{solution}

\begin{solution}{exo:g11:titration:10}
$n(\ce{Fe^{2+}}) = 5 \times 0.0200 \times 0.0130 = \qty{1.30e-3}{mol}$ in
\qty{10.0}{mL}: $c = \qty{0.130}{mol/L}$ voor de verdunde oplossing, en
$10 \times 0.130 = \qty{1.30}{mol/L}$ voor de geconcentreerde.
\end{solution}

\begin{solution}{exo:g11:titration:11}
$0.0630 \times 55.8 = \qty{3.52}{g}$.
\end{solution}

\begin{solution}{exo:g11:titration:12}
$0.1 / 7.2 \approx \qty{1.4}{\percent}$, tegen
\qty{0.7}{\percent} voor \qty{14.3}{mL}. Titreer een groter volume (of
een geconcentreerder monster) van de oplossing, of gebruik een meer
verdunde \omterm{def:g11:titration:titration}{titrant}.
\end{solution}

\begin{solution}{exo:g11:titration:13}
$n(\ce{MnO4^-}) = 0.0200 \times 0.0176 = \qty{3.52e-4}{mol}$;
$n(\ce{H2O2}) = \frac{5}{2} \times \num{3.52e-4} = \qty{8.80e-4}{mol}$ in
\qty{10.0}{mL}: \qty{0.0880}{mol/L} in de verdunde oplossing,
\qty{0.880}{mol/L} in het ontsmettingsmiddel, dat is
$0.880 \times 34.0 = \qty{29.9}{g/L}$ (ongeveer een oplossing van 3 \%).
\end{solution}

\begin{solution}{exo:g11:titration:14}
$n(\ce{MnO4^-}) = \num{1.4e-3} / 5 = \qty{2.8e-4}{mol}$, toe te voegen
in ongeveer \qty{0.015}{L}: $c \approx \qty{0.019}{mol/L}$, dus een
oplossing van \qty{0.0200}{mol/L}. Bij tien keer zo geconcentreerd zou
$V_E$ ongeveer \qty{1.4}{mL} zijn: de fout van de buret zou dan zo'n
\qty{7}{\percent} worden.
\end{solution}

\begin{solution}{exo:g11:titration:15}
$8 \times 0.0200 \times 0.0143 = \qty{2.29e-3}{mol}$ waterstofionen.
\qty{10}{mL} van \qty{1.0}{mol/L} bevat \qty{1.0e-2}{mol}: genoeg, ruim
vier keer zoveel. Zonder zuur zou de reactie van dit hoofdstuk niet
kunnen verlopen zoals ze geschreven is (waterstofionen zijn een
\omterm{def:g7:chemical-reactions:reactant}{beginstof}), en zou permanganaat anders reageren.
\end{solution}

\begin{solution}{pb:g11:titration:1}
\textbf{1.} \ce{MnO4^-}/\ce{Mn^{2+}} en \ce{Fe^{3+}}/\ce{Fe^{2+}}. De
\omterm{def:g11:redox:oxidant}{oxidator} is het permanganaation, de \omterm{def:g11:redox:oxidant}{reductor} het ijzer(II)\omterm{def:g9:ions:ion}{ion}.

\textbf{2.} \ce{MnO4^- + 8H+ + 5e- -> Mn^{2+} + 4H2O};
\ce{Fe^{2+} -> Fe^{3+} + e-}.

\textbf{3.} De tweede wordt met 5 vermenigvuldigd:
\ce{MnO4^- + 8H+ + 5Fe^{2+} -> Mn^{2+} + 5Fe^{3+} + 4H2O}. Lading
links $-1 + 8 + 10 = +17$; rechts $+2 + 15 = +17$.

\textbf{4.} Waterstofionen zijn een \omterm{def:g7:chemical-reactions:reactant}{beginstof}: de reactie heeft een \omterm{def:g9:acids-bases-ph:acidic}{zure
oplossing} nodig.

\textbf{5.} Ze is aflopend, snel en de enige reactie van het
permanganaat in de erlenmeyer; het paarse permanganaat wordt vóór het
\omterm{def:g11:titration:equivalence}{equivalentiepunt} verbruikt en kleurt de erlenmeyer daarna roze.

\textbf{6.} In de buret; aflezen aan de onderkant van de meniscus, op
ooghoogte, vóór en na, tot op \qty{0.05}{mL}.

\textbf{7.} Lichtgeel (ijzerionen); elke paarse druppel wordt meteen
verbruikt door de ijzer(II)\omterm{def:g9:ions:ion}{ionen} die er nog zijn.

\textbf{8.} $0.0200 \times 0.0143 = \qty{2.86e-4}{mol}$.

\textbf{9.} $n(\ce{Fe^{2+}}) = 5 \times \num{2.86e-4} =
\qty{1.43e-3}{mol}$.

\textbf{10.} $\num{1.43e-3} \times 55.8 = \qty{0.0798}{g} =
\qty{79.8}{mg}$.

\textbf{11.} $(80 - 79.8)/80 \approx \qty{0.3}{\percent}$.

\textbf{12.} $\qty{0.1}{mL}$ op \qty{14.3}{mL} is \qty{0.7}{\percent},
ongeveer \qty{0.6}{mg}: het resultaat, $79.8 \pm 0.6$ mg, klopt met het
etiket.

\textbf{13.} $14.1$: $5 \times 0.0200 \times 0.0141 \times 55.8 =
\qty{78.7}{mg}$; $14.4$: \qty{80.4}{mg}.

\textbf{14.} $(79.8 + 78.7 + 80.4)/3 = \qty{79.6}{mg}$.

\textbf{15.} De tabletten zelf verschillen een beetje in ijzergehalte,
ongeveer \qty{1}{\percent}.

\textbf{16.} Nee: $V_E$ zou ongeveer \qty{2.9}{mL} zijn, en de fout van
de buret \qty{3.5}{\percent}, vijf keer slechter.

\textbf{17.} Vitamine C zou ook door het permanganaat geoxideerd worden:
de reactie zou niet langer eenduidig zijn, $V_E$ zou te groot worden, en
de hoeveelheid ijzer zou overschat worden.

\textbf{18.} $\num{1.43e-3} \times 151.9 = \qty{0.217}{g}$, ongeveer
\qty{217}{mg}.

\textbf{19.} $\num{1.43e-3} \times \num{6.02e23} = \num{8.6e20}$ \omterm{def:g9:ions:ion}{ionen}.

\textbf{20.} Ongeveer \qty{80}{mg} ijzer per tablet: het etiket is
eerlijk.
\end{solution}
